\documentclass[11pt]{article}

\usepackage[T1]{fontenc}
\usepackage[utf8]{inputenc}
\usepackage[english]{babel}
\usepackage{amsmath,amssymb,amsthm}
\usepackage{array}
\usepackage{booktabs}
\usepackage{geometry}
\usepackage{placeins}
\usepackage{hyperref}
\usepackage{url}

\geometry{margin=1in}

\hypersetup{
  colorlinks=true,
  linkcolor=black,
  citecolor=blue,
  urlcolor=blue,
  pdftitle={Computational Extension of OEIS Sequence A163665},
  pdfauthor={Carlo Corti}
}

\newtheorem{theorem}{Theorem}
\newtheorem{proposition}[theorem]{Proposition}
\newtheorem{lemma}[theorem]{Lemma}
\newtheorem{corollary}[theorem]{Corollary}
\theoremstyle{definition}
\newtheorem{definition}[theorem]{Definition}
\theoremstyle{remark}
\newtheorem{remark}[theorem]{Remark}

\newcommand{\seqid}{A163665}

\title{Computational Extension of OEIS Sequence A163665:\\
Eight New Terms via a Segmented Prime-Index Search}

\author{Carlo Corti\\
Independent researcher\\
\texttt{carlo.corti@outlook.com}}

\date{July 2026}

\begin{document}

\maketitle

\begin{abstract}
OEIS \seqid{} consists of the integers \(k\) such that
\[
  k=\operatorname{prime}(\varphi(\sigma(\pi(k)))).
\]
Since every term is prime, the search can be transferred to the prime
index \(n=\pi(k)\): an index \(n\) is a hit precisely when
\(\varphi(\sigma(n))=n\), and the corresponding sequence term is
\(\operatorname{prime}(n)\).  This note reports a segmented
ISO/IEC 9899:2018 C (C17) \cite{isoC17}
computation over
\[
  234954224 \leq n \leq 405234954223
\]
the range following the earlier OEIS bound \(k \leq 5\cdot 10^9\).
The campaign found eight new terms and no further terms through
\[
  \operatorname{prime}(405234954223)=11776117514717
\]
The negative certification is finite and computational: it holds subject
to the correctness of the segmented production program and the integrity
of the recorded run manifests.  All 28 positive rows in the consolidated
b-file were also checked independently with PARI/GP for
\(\varphi(\sigma(n))=n\) and primality of \(k\), and with primecount for
the identity
\(k=\operatorname{prime}(n)\); this includes the eight new terms.
\end{abstract}

\noindent\textbf{2020 Mathematics Subject Classification.}
Primary 11Y55; Secondary 11A25, 11Y11.

\noindent\textbf{Key words and phrases.}
integer sequence, OEIS, Euler totient function, sum of divisors,
prime index, computational number theory.

\section{Introduction}

Let \(\operatorname{prime}(m)\) denote the \(m\)-th prime, let \(\pi(x)\)
denote the prime-counting function, let \(\varphi\) be Euler's totient
function, and let \(\sigma\) be the sum-of-divisors function.  OEIS
\seqid{} is defined by
\[
  k=\operatorname{prime}(\varphi(\sigma(\pi(k))))
\]
The OEIS snapshot used for this computation lists the first 20 terms,
ending at
\[
  a(20)=2444540149
\]
and records that there is no further term up to \(5\cdot 10^9\)
\cite{oeisA163665}.  Since
\[
  \pi(5\cdot 10^9)=234954223
\]
the next natural computational task is to search prime indices
\(n>234954223\).

The main contribution of this note is a finite computational extension.
The production campaign scans the contiguous interval
\[
  234954224 \leq n \leq 405234954223
\]
and reports exactly eight hits.  Translating those hits back by
\(k=\operatorname{prime}(n)\) gives
\[
\begin{gathered}
50685770167,\quad 94206075193,\quad 106882008719,\quad
571264143487,\\
1272242879459,\quad 2536961181761,\quad
3715374607207,\quad 9576704442589.
\end{gathered}
\]
The computation also certifies that there is no further term after
\(9576704442589\) up to
\[
  11776117514717
\]

This is a computational note rather than a theoretical classification.
No density law, infinite family, or proof of eventual behavior is claimed.
The finite exclusion interval is certified by the C17 segmented kernel,
the recorded run manifest, and the validation artifacts described below.

\section{Definition and reduction to prime indices}

\begin{definition}
An integer \(k\geq 2\) belongs to \seqid{} if
\[
  k=\operatorname{prime}(\varphi(\sigma(\pi(k))))
\]
\end{definition}

\begin{proposition}[Prime-index reduction]
\label{prop:reduction}
Let \(p_n=\operatorname{prime}(n)\).  An integer \(k\) belongs to
\seqid{} if and only if \(k=p_n\) for some \(n\geq 1\) satisfying
\[
  \varphi(\sigma(n))=n
\]
\end{proposition}

\begin{proof}
If \(k\in\seqid{}\), then the right hand side of the defining equation is
a prime, so \(k\) is prime.  Hence \(k=p_n\), where \(n=\pi(k)\).  The
defining equation becomes
\[
  p_n=\operatorname{prime}(\varphi(\sigma(n)))
\]
and injectivity of the prime-index map gives \(\varphi(\sigma(n))=n\).
Conversely, if \(\varphi(\sigma(n))=n\) and \(k=p_n\), then \(\pi(k)=n\) and
therefore
\[
  \operatorname{prime}(\varphi(\sigma(\pi(k))))
  =\operatorname{prime}(\varphi(\sigma(n)))
  =\operatorname{prime}(n)=k
\]
\end{proof}

\begin{lemma}[Parity filter]
\label{lem:parity}
Every solution \(n>1\) of \(\varphi(\sigma(n))=n\) is even.
\end{lemma}

\begin{proof}
For \(n>1\), one has \(\sigma(n)>2\).  Euler's totient is even for every
integer greater than 2.  Thus \(\varphi(\sigma(n))\) is even, and so \(n\) is
even whenever \(\varphi(\sigma(n))=n\).
\end{proof}

Lemma~\ref{lem:parity} is the only global pruning rule used in the
certified scan.  Stronger congruence restrictions were not assumed.

\section{Computational method}

The production program is the single C17 source file
\texttt{src/a163665.c}, built by \texttt{src/Makefile}.  It provides
three modes:
\[
\texttt{selftest},\qquad
\texttt{segment --lo N --hi M --threads T --out PREFIX},\qquad
\texttt{summarize}.
\]
The search mode treats the range \([N,M]\) as inclusive.  Within each
segment it retains \(n=1\) and even \(n>1\), in accordance with
Lemma~\ref{lem:parity}.

For each retained candidate \(n\), the kernel performs the following
exact computation.
\begin{enumerate}
\item Factor \(n\) within the segment and compute \(\sigma(n)\) from its
factorization.  The source builds a prime table up to \(\sqrt{\texttt{hi}}\),
visits the multiples of each prime in the inclusive segment, removes
prime powers from a mutable residual array, and multiplies the
corresponding divisor-sum factors into a checked 64-bit accumulator.
Any residual factor larger than 1 is then accounted for as a prime
factor of \(n\).
\item Factor \(\sigma(n)\) using the 64-bit factorization routines in the
production source.  These first remove prime factors up to 97, then use a
deterministic Miller--Rabin test with bases
\[
2,\ 325,\ 9375,\ 28178,\ 450775,\ 9780504,\ 1795265022
\]
and a fixed-seed Pollard--Rho recursion for composite cofactors
\cite{miller1976,rabin1980,sprpBases,pollard1975}.  This base set is
recorded as deterministic for at least the unsigned 64-bit range; the
implementation reduces each base modulo the candidate before testing.
The
Pollard--Rho routine tries constants \(c=1,\ldots,32\), uses the
iteration \(x\mapsto x^2+c\pmod n\) with Floyd cycle detection, and
declares a candidate unresolved if factorization fails.  Arithmetic is
checked in unsigned 64-bit storage, with 128-bit intermediates for
products and modular arithmetic; an overflow or factorization failure is
therefore a segment failure or unresolved candidate, not a skipped case.
\item For every prime divisor \(q\) of \(\sigma(n)\), reject the candidate
if \(q-1\nmid n\).  This is a necessary condition, not a heuristic:
if \(\varphi(\sigma(n))=n\), then each factor \(q-1\) occurring in
\(\varphi(\sigma(n))\) divides \(n\).
\item Compute \(\varphi(\sigma(n))\) exactly from the factorization of
\(\sigma(n)\).  If the result equals \(n\), record a hit.
\item Convert a hit index to a sequence term by computing
\(\operatorname{prime}(n)\).  The production source uses an internal
Lehmer prime-counting routine for small hit indices and delegates larger
hit-index conversion to the command-line program \texttt{primecount};
the public validation script independently rechecks every reported
mapping with \texttt{primecount --nth-prime --threads=1}.
\end{enumerate}

Each segment writes a hit table and a manifest.  The manifest records the
range, parity policy, tested-candidate count, rejected-candidate count,
factorization-path telemetry, hit count, status, elapsed time, and file
checksums.  The source and output artifacts are protected by SHA-256
fields; compact FNV-1a 64-bit fields are also recorded as secondary
checksums.  Only segments with status \texttt{OK\_EMPTY} or
\texttt{OK\_HITS}, zero unresolved candidates, and zero failed candidates
are used in the certified negative result.  The manuscript does not claim
formal verification of the C implementation or a second exhaustive scan
of the non-hit interval.

\section{Validation and independent checks}

The public package records the final validation state in
\texttt{validation/validation\_summary.md}.  The final campaign synthesis
uses official runs \texttt{A163665\_run\_001} through
\texttt{A163665\_run\_015}; the pilot overlap run
\texttt{A163665\_run\_000} and interrupted old-binary artifacts are
excluded.

The final validation report gives the following aggregate counts:
\[
\begin{array}{lr}
\text{official runs} & 15\\
\text{certified subsegments} & 40500\\
\text{raw }n\text{-width} & 405000000000\\
\text{candidates tested} & 202500000000\\
\text{unresolved segments} & 0\\
\text{failed segments} & 0\\
\text{new hits} & 8
\end{array}
\]
The elapsed total computed from the official run logs is \(22h{:}12m{:}58s\).
This is the elapsed total reconstructed from the official run-log time
fields, not a mathematical input and not a portable benchmark.  The
public aggregate manifest records run-level coverage, hit counts, and
failure/unresolved status; it does not by itself preserve every
per-subsegment timing and thread-count field.

The local validation environment for this package was Linux Mint 22.3 on
an AMD Ryzen 9 7940HS system with 8 cores, 16 hardware threads, and
64 GB DDR5 installed RAM.  The public Makefile builds with GCC 13.3.0
using
\[
\texttt{-std=c17 -O2 -Wall -Wextra -Wshadow -Wformat=2}
\]
together with the additional warning flags in \texttt{src/Makefile} and
\texttt{-fopenmp}.  The independent scripts were run with PARI/GP 2.15.4
and primecount 7.10 in the local validation environment.  These version
strings document the package validation environment; final publication
should also record the exact Git commit or release tag used for the
archived package.

The validation matrix is:
\begin{center}
\begin{tabular}{>{\raggedright\arraybackslash}p{0.48\textwidth}
>{\raggedright\arraybackslash}p{0.38\textwidth}}
\toprule
\textbf{Check} & \textbf{Recorded result}\\
\midrule
C selftest: b-file terms, known index witnesses, low-range cross-check,
and segmented/direct comparison & \texttt{selftest OK}\\
Global run-summary validation & \texttt{bad=0}, contiguous coverage\\
Official submanifest count & \(40500\) files\\
One-point reproduction of the largest discovered hit
& \texttt{OK\_HITS} at \(n=331914240000\)\\
PARI/GP positive-hit verification
& all 28 terms OK\\
primecount prime-index mapping verification
& all 28 mappings OK\\
\bottomrule
\end{tabular}
\end{center}

The PARI/GP script \path{validation/verify_hits.gp} recomputes
\(\varphi(\sigma(n))=n\) for all 28 consolidated positive indices and checks
primality of the corresponding \(k\)-values.  Thus the independent
positive checker covers the 20 previously listed OEIS terms as well as
the eight new discoveries.  The shell script
\path{validation/verify_prime_mapping.sh} uses primecount to verify
\(k=\operatorname{prime}(n)\) for the same 28 rows.  These scripts are
independent positive-hit checks.  They do not constitute a second
exhaustive scan of the non-hit interval.  The negative certification is
therefore bounded by the correctness of the production kernel and the
integrity of the recorded segment manifests.
The endpoint value used for the final bound,
\[
  \operatorname{prime}(405234954223)=11776117514717
\]
is recorded in \path{validation/validation_summary.md} and is reproduced
by
\[
  \texttt{primecount 405234954223 --nth-prime --threads=1}
\]
The PARI/GP checker intentionally avoids \(\pi(k)\) or
\(\operatorname{prime}(n)\) calls at the largest indices, because those
are not lightweight checks in PARI/GP; the prime-index mapping is checked
separately with primecount.

\section{Results}

\begin{theorem}[Validated production search result]
\label{thm:main}
Subject to the correct execution of the validated C17 production kernel
and to the integrity of the recorded official run manifests, the only
solutions of
\[
  \varphi(\sigma(n))=n
\]
in the interval
\[
  234954224 \leq n \leq 405234954223
\]
are
\[
\begin{gathered}
2147483648,\quad 3889036800,\quad 4389396480,\quad
21946982400,\\
47416320000,\quad 92177326080,\quad 133145026560,\quad
331914240000.
\end{gathered}
\]
Their corresponding \seqid{} terms are the eight values listed in
Table~\ref{tab:newterms}.
\end{theorem}

\begin{proof}[Proof, computational]
The interval was partitioned into the 15 official runs listed in
Appendix~\ref{app:runs}.  The run-level manifest
\texttt{validation/run\_manifest.tsv} records contiguous coverage from
\(234954224\) through \(405234954223\), zero unresolved segments, zero
failed segments, and eight total hits.  The final validation report
records \(40500\) certified subsegments and confirms the same global
coverage and hit count.  The hit rows in
\texttt{results/A163665\_campaign\_hits.tsv} and
\texttt{data/certified\_terms.tsv} are identical up to comment/header
lines.  The independent PARI/GP and primecount scripts verify all 28
positive rows in the consolidated b-file, including the eight new rows
reported in Table~\ref{tab:newterms}.
\end{proof}

\begin{table}[ht]
\centering
\begin{tabular}{rrrr}
\toprule
\(a(i)\) & \(n\) & \(k=\operatorname{prime}(n)\) & \(\sigma(n)\)\\
\midrule
21 & 2147483648 & 50685770167 & 4294967295\\
22 & 3889036800 & 94206075193 & 16331433888\\
23 & 4389396480 & 106882008719 & 18377794080\\
24 & 21946982400 & 571264143487 & 94951936080\\
25 & 47416320000 & 1272242879459 & 204721968000\\
26 & 92177326080 & 2536961181761 & 386940247200\\
27 & 133145026560 & 3715374607207 & 601662398400\\
28 & 331914240000 & 9576704442589 & 1433565580920\\
\bottomrule
\end{tabular}
\caption{The eight new certified terms of OEIS \seqid{}.  For each row,
the certified equality is \(\varphi(\sigma(n))=n\).}
\label{tab:newterms}
\end{table}

\begin{corollary}[Extended OEIS bound]
\label{cor:bound}
With the same computational assumptions as in Theorem~\ref{thm:main},
and using the previously recorded OEIS verification that no further term
occurs for \(k\leq 5\cdot 10^9\),
the terms of \seqid{} not exceeding \(11776117514717\) are exactly the
28 terms in the consolidated b-file, ending at
\[
  a(28)=9576704442589
\]
and there is no further term in
\[
  9576704442589 < k \leq 11776117514717
\]
\end{corollary}

\begin{proof}
By Proposition~\ref{prop:reduction}, terms \(k\) of \seqid{} in the
searched range are exactly \(k=\operatorname{prime}(n)\) for the index
hits \(n\) in Theorem~\ref{thm:main}.  The final endpoint is recorded as
\[
  \operatorname{prime}(405234954223)=11776117514717
\]
Since the previous OEIS-style bound corresponds to
\(\pi(5\cdot 10^9)=234954223\), the certified index interval begins
immediately after that bound.
\end{proof}

The file \path{data/b163665.txt} is therefore the candidate extended
OEIS b-file associated with this computation.

\section{Observed structure and conjectural remarks}

The eight new hit indices have relatively smooth factorizations; see
Appendix~\ref{app:factors}.  Several rows contain a common factor
\(2^{13}\), and the factorizations of \(\sigma(n)\) repeatedly involve
small primes such as \(43\) and \(127\).  In these eight new rows, the
largest prime factor of \(n\) is at most \(31\).  The first new hit,
\[
  n=2147483648=2^{31}
\]
is especially transparent:
\[
  \sigma(n)=2^{32}-1
  =3\cdot 5\cdot 17\cdot 257\cdot 65537
\]
and the product of \(q-1\) over the distinct prime factors \(q\) gives
\(\varphi(\sigma(n))=2^{31}\).
Thus the largest prime factor of \(\sigma(n)\) over all eight new rows
is \(65537\), while among the remaining seven rows it is \(2801\).

These features are reported only as observations.  They were not used as
filters in the certified search, and they do not presently constitute an
infinite-family construction or a density model.  A cautious conjectural
reading is that future hits, if any, may continue to favor indices whose
factorization makes \(\sigma(n)\) highly structured.  No claim stronger
than this is supported by the present computation.

\section{Reproducibility}

The raw GitHub-ready package contains the production source, Makefile,
data tables, validation reports, and independent scripts required to
audit the finite computation.  The main local artifacts are:
\[
\begin{array}{ll}
\texttt{src/a163665.c} & \text{C17 production source},\\
\texttt{src/Makefile} & \text{build, test, and sanitizer targets},\\
\texttt{docs/confirmation\_A163665.md} & \text{OEIS snapshot excerpt and triage confirmation},\\
\texttt{docs/strategy\_A163665.md} & \text{computational strategy notes},\\
\texttt{validation/run\_manifest.tsv} & \text{official run-level manifest},\\
\texttt{validation/validation\_summary.md} & \text{final validation report},\\
\texttt{data/b163665.txt} & \text{candidate extended OEIS b-file},\\
\texttt{data/certified\_terms.tsv} & \text{new terms with witnesses},\\
\texttt{results/A163665\_campaign\_hits.tsv} & \text{consolidated hit table},\\
\texttt{validation/verify\_hits.gp} & \text{PARI/GP positive-hit checker},\\
\texttt{validation/verify\_prime\_mapping.sh} & \text{primecount mapping checker}.
\end{array}
\]

The package validation report records the following commands as passed:
\[
\begin{array}{l}
\texttt{make -C src clean all}\\
\texttt{make -C src test THREADS=4}\\
\texttt{make -C src sanitize}\\
\texttt{gp -q validation/verify\_hits.gp}\\
\texttt{validation/verify\_prime\_mapping.sh}\\
\texttt{make -C src clean}.
\end{array}
\]
The large local subsegment directories generated during the campaign are
not included in the raw public package.  Their aggregate certified
content is represented by the run manifest and validation summary.
The final public release should pin a Git commit or release tag and
replace the Zenodo placeholder in the bibliography by an actual DOI.

\section{Conclusion}

The computation extends OEIS \seqid{} from 20 to 28 known terms.  The
new terms are the eight primes listed in Table~\ref{tab:newterms}, and
the certified absence of additional terms now reaches
\[
  k=11776117514717
\]
The mathematical reduction to the fixed-point condition
\(\varphi(\sigma(n))=n\) is exact, while the extension itself is a finite
computational result.  The main limitation is also explicit: the positive
rows have independent PARI/GP and primecount checks, but the exclusion of
non-hits over the full interval rests on the validated segmented C17
campaign and its recorded manifests.

\section*{Data availability}

The code, OEIS snapshot excerpt and confirmation notes, result tables,
manifests, validation summaries, candidate b-file, and independent
positive-hit checkers are part of the project package and are intended to
be included in the GitHub repository and archive listed in the References
when those public artifacts are finalized.  The snapshot material is
recorded in \texttt{docs/confirmation\_A163665.md}, not as a separate
standalone input-snapshot file.  Until the public repository tag and
Zenodo DOI are assigned, the manuscript should be read as a pre-release
computational note whose negative certificate is auditable from the local
package but not yet from an immutable public archive.  No additional
dataset is needed to reproduce the finite claims stated in this note,
apart from standard toolchain components and the external programs used
by the recorded validation commands.

\section*{AI use disclosure}

Large language models were used as research and drafting assistants during the
\seqid{} workflow, including OpenAI Codex and external LLM-based review systems
from the ChatGPT, Claude, DeepSeek, Gemini, Grok, Kimi, Qwen, MiniMax and Z
families.  NotebookLM was also used as a research-assistant tool during the
methodological workflow.  These systems are named at family or tool level where
exact model-version metadata were not consistently available in the local
review artifacts.

Project workflow.  LLM assistance was used to organize the OEIS snapshot,
triage the equivalence between the original definition and the prime-index
fixed-point problem, compare possible computational strategies, and prepare
publication-facing notes from the consolidated package.  Internal triage
judgments and convenience assessments are not used as mathematical evidence in
this paper.

Mathematical analysis.  LLM assistance was used to structure the reduction from
\(k\) to \(n=\pi(k)\), the parity filter, the exact \((q-1)\mid n\) rejection
condition, and the distinction between theorem, computation, and conjectural
observation.  The mathematical acceptance criterion is the explicit equality
\(\varphi(\sigma(n))=n\), checked against the local OEIS snapshot, the result tables,
and the independent PARI/GP positive-hit verifier.

Algorithm and code.  LLM assistance was used in the design, explanation, and
review loop for the C17 segmented implementation.  The final numerical claims
rely on the local source, Makefile, run manifest, validation summary, b-file,
known-term replay, sanitizer run, one-point reproduction of the largest new
hit, and independent positive-hit scripts, not on conversational assertions.

Manuscript.  LLM assistance was used to draft and structure this manuscript
from the project artifacts.  Numerical claims in the manuscript were checked
against the local b-file, certified-term table, campaign hit table, run
manifest, validation notes, validation summary, and independent verifier
scripts.

The author takes full scientific responsibility for the mathematical
statements, algorithm, implementation, numerical results, validation
boundaries, and conclusions of this paper.

\clearpage

\appendix

\section{Official run manifest summary}
\label{app:runs}

\begin{table}[ht]
\centering
\begin{tabular}{lrrr}
\toprule
Run & \(n\)-interval & Certified subsegments & Hits\\
\midrule
\texttt{run\_001} & \(234954224\)--\(30234954223\) & 3000 & 4\\
\texttt{run\_002} & \(30234954224\)--\(60234954223\) & 3000 & 1\\
\texttt{run\_003} & \(60234954224\)--\(90234954223\) & 3000 & 0\\
\texttt{run\_004} & \(90234954224\)--\(120234954223\) & 3000 & 1\\
\texttt{run\_005} & \(120234954224\)--\(150234954223\) & 3000 & 1\\
\texttt{run\_006} & \(150234954224\)--\(180234954223\) & 3000 & 0\\
\texttt{run\_007} & \(180234954224\)--\(205234954223\) & 2500 & 0\\
\texttt{run\_008} & \(205234954224\)--\(230234954223\) & 2500 & 0\\
\texttt{run\_009} & \(230234954224\)--\(255234954223\) & 2500 & 0\\
\texttt{run\_010} & \(255234954224\)--\(280234954223\) & 2500 & 0\\
\texttt{run\_011} & \(280234954224\)--\(305234954223\) & 2500 & 0\\
\texttt{run\_012} & \(305234954224\)--\(330234954223\) & 2500 & 0\\
\texttt{run\_013} & \(330234954224\)--\(355234954223\) & 2500 & 1\\
\texttt{run\_014} & \(355234954224\)--\(380234954223\) & 2500 & 0\\
\texttt{run\_015} & \(380234954224\)--\(405234954223\) & 2500 & 0\\
\bottomrule
\end{tabular}
\caption{Official run-level manifest summary.}
\label{tab:runmanifest}
\end{table}

\FloatBarrier

Every official run in the public manifest has
\texttt{coverage=CONTIGUOUS}, \texttt{segments\_unresolved=0}, and
\texttt{segments\_failed=0}.

\section{Factorization witnesses for the new hits}
\label{app:factors}

\scriptsize
\begin{table}[ht]
\centering
\begin{tabular}{r>{\raggedright\arraybackslash}p{0.31\textwidth}
>{\raggedright\arraybackslash}p{0.43\textwidth}}
\toprule
\(n\) & Factorization of \(n\) & Factorization of \(\sigma(n)\)\\
\midrule
2147483648 &
\(2^{31}\) &
\(3\cdot 5\cdot 17\cdot 257\cdot 65537\)\\
3889036800 &
\(2^9\cdot 3^4\cdot 5^2\cdot 11^2\cdot 31\) &
\(2^5\cdot 3\cdot 7\cdot 11^3\cdot 19\cdot 31^2\)\\
4389396480 &
\(2^{13}\cdot 3^7\cdot 5\cdot 7^2\) &
\(2^5\cdot 3^3\cdot 5\cdot 19\cdot 41\cdot 43\cdot 127\)\\
21946982400 &
\(2^{13}\cdot 3^7\cdot 5^2\cdot 7^2\) &
\(2^4\cdot 3^2\cdot 5\cdot 19\cdot 31\cdot 41\cdot 43\cdot 127\)\\
47416320000 &
\(2^{13}\cdot 3^3\cdot 5^4\cdot 7^3\) &
\(2^7\cdot 3\cdot 5^3\cdot 11\cdot 43\cdot 71\cdot 127\)\\
92177326080 &
\(2^{13}\cdot 3^8\cdot 5\cdot 7^3\) &
\(2^5\cdot 3^2\cdot 5^2\cdot 13\cdot 43\cdot 127\cdot 757\)\\
133145026560 &
\(2^{13}\cdot 3^6\cdot 5\cdot 7^3\cdot 13\) &
\(2^6\cdot 3^2\cdot 5^2\cdot 7\cdot 43\cdot 127\cdot 1093\)\\
331914240000 &
\(2^{13}\cdot 3^3\cdot 5^4\cdot 7^4\) &
\(2^3\cdot 3\cdot 5\cdot 11\cdot 43\cdot 71\cdot 127\cdot 2801\)\\
\bottomrule
\end{tabular}
\caption{Factorization witnesses for the eight new hits.}
\label{tab:factors}
\end{table}
\normalsize

\clearpage

\begin{thebibliography}{99}
\sloppy
\raggedright

\bibitem{oeisA163665}
OEIS Foundation Inc.,
\emph{Sequence A163665: Numbers \(k\) such that
\(k=\operatorname{prime}(\varphi(\sigma(\pi(k))))\)},
The On-Line Encyclopedia of Integer Sequences,
\url{https://oeis.org/A163665}.
Submitted by Farideh Firoozbakht, Aug. 4, 2009.  Local project snapshot
used as source material, July 6, 2026.

\bibitem{cortiA163665Package2026}
C. Corti,
\emph{A163665: C17 segmented search and validation package},
GitHub repository and Zenodo archive, 2026.
\url{https://github.com/carcorti/A163665}.
Zenodo DOI placeholder: \texttt{10.5281/zenodo.xxxxxxxx}, to be
replaced after archive publication.

\bibitem{miller1976}
G. L. Miller,
\emph{Riemann's hypothesis and tests for primality},
Journal of Computer and System Sciences 13 (1976), 300--317.

\bibitem{rabin1980}
M. O. Rabin,
\emph{Probabilistic algorithm for testing primality},
Journal of Number Theory 12 (1980), 128--138.

\bibitem{sprpBases}
W. Izykowski,
\emph{Deterministic variants of the Miller--Rabin primality test:
SPRP bases}.
\url{https://miller-rabin.appspot.com/}.
Used for the recorded seven-base strong probable-prime set covering at
least \(2^{64}\).

\bibitem{pollard1975}
J. M. Pollard,
\emph{A Monte Carlo method for factorization},
BIT Numerical Mathematics 15 (1975), 331--334.

\bibitem{isoC17}
International Organization for Standardization,
\emph{ISO/IEC 9899:2018, Information technology -- Programming languages -- C},
2018.

\bibitem{pariGP}
The PARI Group,
\emph{PARI/GP computer algebra system}, version 2.15.4.
\url{https://pari.math.u-bordeaux.fr/}.
Used here through the project script
\texttt{validation/verify\_hits.gp}.

\bibitem{primecount}
K. Walisch,
\emph{primecount: fast prime counting function library and command-line
program}, version 7.10.
\url{https://github.com/kimwalisch/primecount}.
Used here through the project script
\texttt{validation/verify\_prime\_mapping.sh}.

\end{thebibliography}

\end{document}
